% CORRECTED VERSION
% Changes:
% - Original proof was missing due to gateway timeout; generated a complete proof from scratch.
% - Added all necessary steps, assumptions, and derivations for convergence of gradient descent on L-smooth convex functions.
% - Ensured mathematical rigor, explicit assumption usage, and step-by-step annotations.

\documentclass{article}
\usepackage{amsmath, amssymb, amsthm, mathtools}
\usepackage{hyperref}

\newtheorem{assumption}{Assumption}
\newtheorem{theorem}{Theorem}
\newtheorem{lemma}{Lemma}
\newtheorem{corollary}{Corollary}
\newtheorem{remark}{Remark}

\title{Convergence Proof for Gradient Descent on Smooth Convex Functions}
\author{}
\date{}

\begin{document}

\maketitle

\section{Problem Setting and Assumptions}

We consider the unconstrained optimization problem:
\begin{equation}
    \min_{x \in \mathbb{R}^d} f(x),
\end{equation}
where $f: \mathbb{R}^d \to \mathbb{R}$ satisfies the following assumptions.

\begin{assumption}[L-smoothness] \label{assump:smooth}
    $f$ is differentiable and its gradient is $L$-Lipschitz continuous for some $L > 0$:
    \begin{equation}
        \|\nabla f(x) - \nabla f(y)\| \le L \|x - y\|, \quad \forall x, y \in \mathbb{R}^d.
    \end{equation}
\end{assumption}

\begin{assumption}[Convexity] \label{assump:convex}
    $f$ is convex:
    \begin{equation}
        f(y) \ge f(x) + \nabla f(x)^\top (y - x), \quad \forall x, y \in \mathbb{R}^d.
    \end{equation}
\end{assumption}

\begin{assumption}[Existence of minimizer] \label{assump:minimizer}
    The set of minimizers $X^* = \arg\min_{x \in \mathbb{R}^d} f(x)$ is nonempty. Let $f^* = \min_x f(x)$ and choose any $x^* \in X^*$.
\end{assumption}

\section{Algorithm}

Gradient Descent (GD) with fixed step size $\eta > 0$:
\begin{equation}
    x_{k+1} = x_k - \eta \nabla f(x_k), \quad k = 0, 1, 2, \dots
\end{equation}

\section{Main Theorem}

\begin{theorem}[Convergence Rate of GD] \label{thm:gd_convergence}
    Under Assumptions \ref{assump:smooth}, \ref{assump:convex}, and \ref{assump:minimizer}, if the step size satisfies $0 < \eta \le \frac{1}{L}$, then for all $k \ge 0$:
    \begin{equation}
        f(x_k) - f^* \le \frac{\|x_0 - x^*\|^2}{2\eta k}.
    \end{equation}
    In particular, with $\eta = \frac{1}{L}$, we have $f(x_k) - f^* \le \frac{L \|x_0 - x^*\|^2}{2k}$.
\end{theorem}

\section{Proof}

The proof proceeds through a series of lemmas and inequalities, each annotated with a step number for traceability.

\begin{lemma}[Descent Lemma] \label{lem:descent}
    Under Assumption \ref{assump:smooth}, for any $x, y \in \mathbb{R}^d$:
    \begin{equation}
        f(y) \le f(x) + \nabla f(x)^\top (y - x) + \frac{L}{2} \|y - x\|^2.
    \end{equation}
    \begin{proof}
        This is a standard consequence of $L$-smoothness (see, e.g., Nesterov, 2004). \qed
    \end{proof}
\end{lemma}

\begin{lemma}[Co-coercivity] \label{lem:cocoercive}
    Under Assumptions \ref{assump:smooth} and \ref{assump:convex}, for any $x, y \in \mathbb{R}^d$:
    \begin{equation}
        (\nabla f(x) - \nabla f(y))^\top (x - y) \ge \frac{1}{L} \|\nabla f(x) - \nabla f(y)\|^2.
    \end{equation}
    \begin{proof}
        Follows from the Baillon-Haddad theorem for convex and $L$-smooth functions. \qed
    \end{proof}
\end{lemma}

Now we prove the main theorem.

\textbf{Step 1:} Apply the Descent Lemma (Lemma \ref{lem:descent}) with $x = x_k$ and $y = x_{k+1} = x_k - \eta \nabla f(x_k)$:
\begin{equation} \label{eq:descent_step}
    f(x_{k+1}) \le f(x_k) - \eta \|\nabla f(x_k)\|^2 + \frac{L \eta^2}{2} \|\nabla f(x_k)\|^2.
\end{equation}
\textbf{Step 2:} Since $\eta \le 1/L$, we have $1 - \frac{L\eta}{2} \ge \frac{1}{2}$. Thus,
\begin{equation} \label{eq:descent_simplified}
    f(x_{k+1}) \le f(x_k) - \frac{\eta}{2} \|\nabla f(x_k)\|^2.
\end{equation}
\textbf{Step 3:} By convexity (Assumption \ref{assump:convex}) with $y = x^*$:
\begin{equation} \label{eq:convexity}
    f(x^*) \ge f(x_k) + \nabla f(x_k)^\top (x^* - x_k).
\end{equation}
Rearranging gives
\begin{equation}
    \nabla f(x_k)^\top (x_k - x^*) \ge f(x_k) - f^*.
\end{equation}
\textbf{Step 4:} Consider the distance to the optimal point:
\begin{align}
    \|x_{k+1} - x^*\|^2 &= \|x_k - \eta \nabla f(x_k) - x^*\|^2 \\
    &= \|x_k - x^*\|^2 - 2\eta \nabla f(x_k)^\top (x_k - x^*) + \eta^2 \|\nabla f(x_k)\|^2. \label{eq:distance_expansion}
\end{align}
\textbf{Step 5:} Using the inequality from Step 3 in \eqref{eq:distance_expansion}:
\begin{equation}
    \|x_{k+1} - x^*\|^2 \le \|x_k - x^*\|^2 - 2\eta (f(x_k) - f^*) + \eta^2 \|\nabla f(x_k)\|^2.
\end{equation}
\textbf{Step 6:} From Step 2, we have $\eta \|\nabla f(x_k)\|^2 \le 2(f(x_k) - f(x_{k+1}))$. Substitute into the previous inequality:
\begin{align}
    \|x_{k+1} - x^*\|^2 &\le \|x_k - x^*\|^2 - 2\eta (f(x_k) - f^*) + 2\eta (f(x_k) - f(x_{k+1})) \\
    &= \|x_k - x^*\|^2 - 2\eta (f(x_{k+1}) - f^*).
\end{align}
\textbf{Step 7:} Rearranging gives
\begin{equation}
    f(x_{k+1}) - f^* \le \frac{1}{2\eta} \left( \|x_k - x^*\|^2 - \|x_{k+1} - x^*\|^2 \right).
\end{equation}
\textbf{Step 8:} Summing this inequality for $k = 0, 1, \dots, K-1$ telescopes the right-hand side:
\begin{equation}
    \sum_{k=0}^{K-1} (f(x_{k+1}) - f^*) \le \frac{1}{2\eta} \|x_0 - x^*\|^2.
\end{equation}
\textbf{Step 9:} Since $f(x_k)$ is non-increasing (by Step 2), we have $f(x_K) \le f(x_{k+1})$ for all $k < K$. Thus,
\begin{equation}
    K (f(x_K) - f^*) \le \sum_{k=0}^{K-1} (f(x_{k+1}) - f^*) \le \frac{1}{2\eta} \|x_0 - x^*\|^2.
\end{equation}
\textbf{Step 10:} Dividing by $K$ yields the desired result:
\begin{equation}
    f(x_K) - f^* \le \frac{\|x_0 - x^*\|^2}{2\eta K}.
\end{equation}
Setting $\eta = 1/L$ gives $f(x_K) - f^* \le \frac{L \|x_0 - x^*\|^2}{2K}$. \qed

\section{Remarks}

\begin{remark}
    The rate $O(1/k)$ is tight for the class of convex $L$-smooth functions under first-order methods.
\end{remark}

\begin{remark}
    If $f$ is additionally $\mu$-strongly convex, a linear convergence rate can be obtained.
\end{remark}

% ===== CORRECTION SUMMARY =====
% 1. Original proof was completely missing (gateway timeout). Generated a full proof from scratch.
% 2. Added all necessary assumptions (A1, A2, A3) and used them explicitly in each step.
% 3. Provided detailed step-by-step derivation with annotations [Step 1] through [Step 10].
% 4. Ensured convergence rate in theorem statement matches the proved rate (O(1/k)).
% 5. Included standard lemmas (Descent Lemma, Co-coercivity) with references for completeness.

\end{document}